% GENERATED FILE. Edit canonical lessons, then run npm run generate:books.
% canonical-source-hash: fnv1a64:1b6c1aa19f533eed
% canonical-lessons: TE-C410-seva, TE-C410-viralam, TE-C410-svacchanda, TE-C410-karyakarta, TE-C410-data

\chapter{Giving and Serving}
\label{ch:te-giving-and-serving}
\hlchaptermodality{410}

\begin{chapteropening}
\textbf{By the end of this chapter:} \emph{I can talk about service, donations and volunteers.}

Complete the last lesson of chapter 410: I can talk about service, donations and volunteers.
\end{chapteropening}

\section[\texorpdfstring{\te{సేవ} (sēva)}{సేవ (sēva)}]{\te{సేవ} (sēva) — service; selfless work}

\label{lesson:TE-C410-seva}



\begin{quote}
Before the new one: say the Telugu for solidarity, then the Telugu for to put into practice.
\end{quote}

\subsection*{You'll want to know: \te{సేవ}}
\textbf{\te{సేవ}} — \emph{sēva} — \textquotedblleft{}service; selfless work\textquotedblright{}.

\begin{quote}
\te{ఆమె} \te{పేదవాళ్ళకు} \te{సేవ} \te{చేస్తుంది}.
\end{quote}

\begin{quote}
\emph{āme pēdavāḷḷaku sēva cēstundi.}
\end{quote}

\begin{quote}
She serves the poor.
\end{quote}

\te{సేవ} \te{చేయు}, to serve; \te{ప్రజా} \te{సేవ}, public service.

\begin{cousinweb}
\textbf{\te{సేవ}} (\emph{sēva}) is borrowed from Sanskrit \emph{sevā}, \textquotedblleft{}going to; service, attendance; devotion\textquotedblright{} (Monier-Williams).
\end{cousinweb}

\begin{grammarlens}[title={What you've built}]
One more word for giving reasons and opinions.
\end{grammarlens}

\subsection*{Your turn}
\begin{itemize}
\raggedright
  \item \emph{Say it:} \emph{sēva}
  \item \emph{Say it:} \emph{sēva}, once more
  \item \emph{Say it:} \emph{āme pēdavāḷḷaku sēva cēstundi.} — She serves the poor.
  \item \emph{Recall:} say the Telugu for solidarity, then the Telugu for to put into practice, then say \emph{sēva} again
\end{itemize}

\begin{tcolorbox}[breakable,colback=teal!4,colframe=teal!35!black,title={Before you move on}]
What does \te{సేవ} mean? (\textquotedblleft{}Service; selfless work\textquotedblright{}.) Say it once more.
\end{tcolorbox}
\section[\texorpdfstring{\te{విరాళం} (virāḷaṁ)}{విరాళం (virāḷaṁ)}]{\te{విరాళం} (virāḷaṁ) — a donation}

\label{lesson:TE-C410-viralam}



\begin{quote}
Before the new one: say the Telugu for to put into practice, then the Telugu for service.
\end{quote}

\subsection*{You'll want to know: \te{విరాళం}}
\textbf{\te{విరాళం}} — \emph{virāḷaṁ} — \textquotedblleft{}a donation\textquotedblright{}.

\begin{quote}
\te{అతను} \te{ఆసుపత్రికి} \te{పెద్ద} \te{విరాళం} \te{ఇచ్చాడు}.
\end{quote}

\begin{quote}
\emph{atanu āsupatriki pedda virāḷaṁ iccāḍu.}
\end{quote}

\begin{quote}
He gave a big donation to the hospital.
\end{quote}

Plural: \te{విరాళాలు}, \emph{virāḷālu}.

\te{విరాళం} \te{ఇచ్చు}, to donate.

\begin{cousinweb}
\textbf{\te{విరాళం}} (\emph{virāḷaṁ}) is a Telugu word whose origin is not settled; neither DEDR nor CDIAL lists it.
\end{cousinweb}

\begin{grammarlens}[title={What you've built}]
One more word for giving reasons and opinions.
\end{grammarlens}

\subsection*{Your turn}
\begin{itemize}
\raggedright
  \item \emph{Say it:} \emph{virāḷaṁ}
  \item \emph{Say it:} \emph{virāḷaṁ}, once more
  \item \emph{Say it:} \emph{atanu āsupatriki pedda virāḷaṁ iccāḍu.} — He gave a big donation to the hospital.
  \item \emph{Recall:} say the Telugu for to put into practice, then the Telugu for service, then say \emph{virāḷaṁ} again
\end{itemize}

\begin{tcolorbox}[breakable,colback=teal!4,colframe=teal!35!black,title={Before you move on}]
What does \te{విరాళం} mean? (\textquotedblleft{}A donation\textquotedblright{}.) Say it once more.
\end{tcolorbox}
\section[\texorpdfstring{\te{స్వచ్ఛంద} (svacchanda)}{స్వచ్ఛంద (svacchanda)}]{\te{స్వచ్ఛంద} (svacchanda) — voluntary}

\label{lesson:TE-C410-svacchanda}



\begin{quote}
Before the new one: say the Telugu for service, then the Telugu for a donation.
\end{quote}

\subsection*{You'll want to know: \te{స్వచ్ఛంద}}
\textbf{\te{స్వచ్ఛంద}} — \emph{svacchanda} — \textquotedblleft{}voluntary\textquotedblright{}.

\begin{quote}
\te{ఆమె} \te{ఒక} \te{స్వచ్ఛంద} \te{సంస్థలో} \te{పని} \te{చేస్తుంది}.
\end{quote}

\begin{quote}
\emph{āme oka svacchanda saṁsthalō pani cēstundi.}
\end{quote}

\begin{quote}
She works for a voluntary organisation.
\end{quote}

\te{స్వచ్ఛంద} \te{సంస్థ}, a voluntary organisation; \te{స్వచ్ఛందంగా}, of one’s own will.

\begin{cousinweb}
\textbf{\te{స్వచ్ఛంద}} (\emph{svacchanda}) is borrowed from Sanskrit \emph{svacchanda}, \textquotedblleft{}one’s own free will; acting at pleasure, independent\textquotedblright{} (Monier-Williams).
\end{cousinweb}

\begin{grammarlens}[title={What you've built}]
One more word for giving reasons and opinions.
\end{grammarlens}

\subsection*{Your turn}
\begin{itemize}
\raggedright
  \item \emph{Say it:} \emph{svacchanda}
  \item \emph{Say it:} \emph{svacchanda}, once more
  \item \emph{Say it:} \emph{āme oka svacchanda saṁsthalō pani cēstundi.} — She works for a voluntary organisation.
  \item \emph{Recall:} say the Telugu for service, then the Telugu for a donation, then say \emph{svacchanda} again
\end{itemize}

\begin{tcolorbox}[breakable,colback=teal!4,colframe=teal!35!black,title={Before you move on}]
What does \te{స్వచ్ఛంద} mean? (\textquotedblleft{}Voluntary\textquotedblright{}.) Say it once more.
\end{tcolorbox}
\section[\texorpdfstring{\te{కార్యకర్త} (kāryakarta)}{కార్యకర్త (kāryakarta)}]{\te{కార్యకర్త} (kāryakarta) — an activist; a volunteer}

\label{lesson:TE-C410-karyakarta}



\begin{quote}
Before the new one: say the Telugu for a donation, then the Telugu for voluntary.
\end{quote}

\subsection*{You'll want to know: \te{కార్యకర్త}}
\textbf{\te{కార్యకర్త}} — \emph{kāryakarta} — \textquotedblleft{}an activist; a volunteer\textquotedblright{}.

\begin{quote}
\te{కార్యకర్తలు} \te{ఇంటింటికీ} \te{వెళ్ళి} \te{సహాయం} \te{చేశారు}.
\end{quote}

\begin{quote}
\emph{kāryakartalu iṇṭiṇṭikī veḷḷi sahāyaṁ cēśāru.}
\end{quote}

\begin{quote}
The volunteers went from house to house and helped.
\end{quote}

Plural: \te{కార్యకర్తలు}, \emph{kāryakartalu}.

A party worker or a volunteer.

\begin{cousinweb}
\textbf{\te{కార్యకర్త}} (\emph{kāryakarta}) is built from Sanskrit \emph{kārya}, \textquotedblleft{}to be done; business\textquotedblright{} (CDIAL 3078), and \emph{kartṛ}, \textquotedblleft{}a doer, an agent\textquotedblright{} (Monier-Williams).
\end{cousinweb}

\begin{grammarlens}[title={What you've built}]
One more word for giving reasons and opinions.
\end{grammarlens}

\subsection*{Your turn}
\begin{itemize}
\raggedright
  \item \emph{Say it:} \emph{kāryakarta}
  \item \emph{Say it:} \emph{kāryakarta}, once more
  \item \emph{Say it:} \emph{kāryakartalu iṇṭiṇṭikī veḷḷi sahāyaṁ cēśāru.} — The volunteers went from house to house and helped.
  \item \emph{Recall:} say the Telugu for a donation, then the Telugu for voluntary, then say \emph{kāryakarta} again
\end{itemize}

\begin{tcolorbox}[breakable,colback=teal!4,colframe=teal!35!black,title={Before you move on}]
What does \te{కార్యకర్త} mean? (\textquotedblleft{}An activist; a volunteer\textquotedblright{}.) Say it once more.
\end{tcolorbox}
\section[\texorpdfstring{\te{దాత} (dāta)}{దాత (dāta)}]{\te{దాత} (dāta) — a donor}

\label{lesson:TE-C410-data}



\begin{quote}
Before the new one: say the Telugu for voluntary, then the Telugu for an activist.
\end{quote}

\subsection*{You'll want to know: \te{దాత}}
\textbf{\te{దాత}} — \emph{dāta} — \textquotedblleft{}a donor\textquotedblright{}.

\begin{quote}
\te{ఒక} \te{దాత} \te{పాఠశాలకు} \te{పుస్తకాలు} \te{ఇచ్చాడు}.
\end{quote}

\begin{quote}
\emph{oka dāta pāṭhaśālaku pustakālu iccāḍu.}
\end{quote}

\begin{quote}
A donor gave books to the school.
\end{quote}

Plural: \te{దాతలు}, \emph{dātalu}.

\te{రక్తదాత}, a blood donor.

\begin{cousinweb}
\textbf{\te{దాత}} (\emph{dāta}) is borrowed from Sanskrit \emph{dātṛ}, \textquotedblleft{}a giver\textquotedblright{} (Monier-Williams), from \emph{dā}, \textquotedblleft{}to give\textquotedblright{}.
\end{cousinweb}

\begin{grammarlens}[title={What you've built}]
One more word for giving reasons and opinions.
\end{grammarlens}

\subsection*{Your turn}
\begin{itemize}
\raggedright
  \item \emph{Say it:} \emph{dāta}
  \item \emph{Say it:} \emph{dāta}, once more
  \item \emph{Say it:} \emph{oka dāta pāṭhaśālaku pustakālu iccāḍu.} — A donor gave books to the school.
  \item \emph{Recall:} say the Telugu for voluntary, then the Telugu for an activist, then say \emph{dāta} again
\end{itemize}

\begin{tcolorbox}[breakable,colback=teal!4,colframe=teal!35!black,title={Before you move on}]
What does \te{దాత} mean? (\textquotedblleft{}A donor\textquotedblright{}.) Say it once more.
\end{tcolorbox}
